\documentclass[11pt]{article}
\usepackage[T1]{fontenc}
\usepackage[utf8]{inputenc}
\usepackage{amsmath,amssymb,amsthm,mathtools}
\usepackage{geometry,graphicx,booktabs,longtable,tabularx,array}
\usepackage{enumitem,needspace,xcolor,tikz,tcolorbox}
\usetikzlibrary{arrows.meta,positioning,calc,fit,backgrounds,decorations.pathreplacing,matrix,patterns}
\tcbuselibrary{breakable}
\usepackage{caption,hyperref,bookmark,cleveref,microtype}
\usepackage{komorebi-publication}
\renewcommand{\KomorebiTitle}{Komorebi Space}
\renewcommand{\KomorebiSubtitle}{Coherent realizations, obstructions, and descent}
\renewcommand{\KomorebiSeries}{Foundations of the Realization Calculus}
\renewcommand{\KomorebiKind}{Monograph}
\renewcommand{\KomorebiEdition}{Monograph edition 1.0.0 / 10 October 2026}
\renewcommand{\KomorebiShortTitle}{Komorebi Space}
\renewcommand{\KomorebiProfile}{monograph}
\KomorebiSetup
\setcounter{tocdepth}{1}
\setcounter{secnumdepth}{2}
\setlist{nosep,leftmargin=1.7em,topsep=4pt}
\hypersetup{pdfkeywords={Komorebi Space,realization calculus,descent,obstruction,Lean,holonomy,closed range}}
\newtheorem{theorem}{Theorem}[section]
\newtheorem{proposition}[theorem]{Proposition}
\newtheorem{lemma}[theorem]{Lemma}
\newtheorem{corollary}[theorem]{Corollary}
\theoremstyle{definition}
\newtheorem{definition}[theorem]{Definition}
\newtheorem{example}[theorem]{Example}
\theoremstyle{remark}
\newtheorem{remark}[theorem]{Remark}
\numberwithin{equation}{section}
\newcommand{\R}{\mathbb R}
\newcommand{\Z}{\mathbb Z}
\newcommand{\N}{\mathbb N}
\newcommand{\K}{\mathbb K}
\newcommand{\V}{\mathcal V}
\newcommand{\W}{\mathcal W}
\newcommand{\id}{\operatorname{id}}
\newcommand{\im}{\operatorname{im}}
\newcommand{\ran}{\operatorname{ran}}
\newcommand{\Hom}{\operatorname{Hom}}
\newcommand{\Ob}{\operatorname{Ob}}
\newcommand{\Ext}{\operatorname{Ext}}
\newcommand{\Match}{\operatorname{Match}}
\newcommand{\Exact}{\operatorname{Exact}}
\newcommand{\dist}{\operatorname{dist}}
\newcommand{\Aut}{\operatorname{Aut}}
\newcommand{\Stab}{\operatorname{Stab}}
\newcommand{\Set}{\mathbf{Set}}
\newcommand{\Cat}{\mathbf{Cat}}
\newcommand{\norm}[1]{\left\lVert#1\right\rVert}
\newcommand{\abs}[1]{\left\lvert#1\right\rvert}
\newcommand{\source}[1]{\par\phantomsection\label{src:#1}{\footnotesize\sffamily\color{KomorebiMuted}Source concordance: \hyperref[app:concordance]{#1}.}\par\smallskip}
\newcommand{\code}[1]{\texttt{\detokenize{#1}}}
\newcommand{\chap}[1]{\par\Needspace{10\baselineskip}\section{#1}}
\newcommand{\partnote}[2]{\clearpage\part{#1}\noindent\begin{minipage}{.92\linewidth}\large #2\end{minipage}\par\vspace{12pt}}
\newtcolorbox{insight}{breakable,colback=KomorebiSoft!55,colframe=KomorebiPine!50,boxrule=.45pt,arc=1pt,left=9pt,right=9pt,top=7pt,bottom=7pt}
\newtcolorbox{scopebox}{breakable,colback=KomorebiWarm,colframe=KomorebiOchre!55,boxrule=.4pt,arc=1pt,left=9pt,right=9pt,top=7pt,bottom=7pt}
\newcommand{\plate}[3]{\begin{figure}[htbp]\centering\input{figures/#1}\caption{#2}\label{#3}\end{figure}}
\tikzset{box/.style={draw=KomorebiPine!70,rounded corners=2pt,fill=KomorebiSoft!50,align=center,inner sep=8pt,font=\small},arr/.style={-{Latex[length=2.2mm]},draw=KomorebiPine,thick},thinarr/.style={-{Latex[length=1.8mm]},draw=KomorebiMuted},lab/.style={font=\small,align=center},dot/.style={circle,fill=KomorebiPine,inner sep=2.3pt}}
\renewcommand{\arraystretch}{1.18}
\raggedbottom
\widowpenalty=7000\clubpenalty=7000
\setlength{\emergencystretch}{3em}

% Compact only the contents entries, without reducing body type size.
\usepackage{etoolbox}
\makeatletter
\patchcmd{\l@section}{1.0em}{0.6em}{}{}
\makeatother

\begin{document}
\KomorebiCover
\pagenumbering{roman}
\section*{About this edition}
This monograph develops Komorebi Space as a theory of coherent realizations relative to a specified system of mathematical views. Its central constructions are ordinary limits, relative extension spaces, obstruction quotients, cocycle groupoids, an explicit stack on a composition site, and quantitative repair operators. The point of the synthesis is to keep these constructions connected without erasing the differences between their hypotheses.

The exposition follows the mathematics rather than the order in which the source library grew. Definitions precede existence questions. Existence is separated from uniqueness. Algebraic feasibility is separated from numerical stability. Changing a representation is separated from changing the data that the representation is required to preserve.

The publication version is \textbf{1.0.0}. This number describes the monograph, not a new universal theorem or a new cumulative Lean release. The accompanying package preserves a complete recovered foundation archive at version 0.20.0 and a separate focused archive, \emph{Moving Exact Spaces}, at version 0.32.0. An additional recovered Hilbert holonomy source is identified separately. The source concordance records where each principal statement is implemented, where several implementations support a combined exposition, and where a printed derivation is an explanatory addition.

Historical compiler and axiom reports retain their original dates and verification strategies. This editorial release performs archive integrity checks, source correspondence checks, exact finite computations, and publication preflight. It does \emph{not} claim a fresh cumulative Lean compilation. The two recovered source trees have not been silently merged into a newly certified library.

All mathematical figures are reproducible vector diagrams or plots of explicit formulas. The source package contains the manuscript, figure sources, exact checks, evidence archives, and the progress ledger. The separate visual guide offers a shorter route through the examples. The Lean guide explains the correspondence and its verification boundaries.

\begin{scopebox}
\textbf{Reading scope.} The word \emph{space} does not assert one universal topological object containing all of mathematics. The strict construction is a set of coherent families. Symmetry is retained in a groupoid, and effective descent is proved for a specified site. Analytic estimates apply only after a normed or Hilbert model has been declared. These are related layers, not interchangeable definitions.
\end{scopebox}

\clearpage
\section*{Orientation}
A description at one scale often omits information visible at another. A formula in one representation may fail to commute with a change of representation. Local pieces may each be valid and still fail to assemble. Even when assembly is possible, there may be many completions, and small discrepancies may or may not imply proximity to an exact answer.

Komorebi Space treats these questions in a common order. First specify the views and the maps between them. Then define the coherent realizations. Next identify the information that must remain fixed. Finally determine the obstruction to extension, the freedom in a successful extension, and the conditions under which repair is stable.

The resulting program has three central distinctions:
\[
\begin{gathered}
\text{a collection of appearances}\quad\ne\quad\text{a coherent realization},\\
\text{existence of an exact answer}\quad\ne\quad\text{a stable repair estimate},\\
\text{one isomorphism class}\quad\ne\quad\text{absence of symmetries}.
\end{gathered}
\]
The monograph makes these distinctions into explicit constructions and theorems.

\subsection*{The route through the book}
Part I defines coherent sections and proves the additive extension criterion. Part II connects the construction to resolution towers and separates finite reconstruction from convergence. Part III retains nonabelian symmetries and proves effective descent on the composition site. Part IV derives geometric repair bounds and controls the motion of exact spaces. The final section states what the framework does and does not establish for mathematical semantics and machine reasoning.

The fastest technical route is Sections~\ref{sec:views}, \ref{sec:obstruction}, \ref{sec:stack}, \ref{sec:closed}, and~\ref{sec:motion}. A reader interested in the role of packets and residuals should include Sections~\ref{sec:split} and~\ref{sec:cascade}. The graph models in Section~\ref{sec:graph} explain how a categorical realization problem acquires an actual numerical condition number.

\subsection*{What is being contributed}
The limit universal property, gauge action groupoids, standard stack definitions, orthogonal projection, and bounded inverse principles belong to established mathematics. They are used here as such. The contribution of this presentation is their explicit integration with the project's actual view systems, protected correction spaces, resolution differentials, graph transports, and checked examples. Historical novelty of every component is not inferred from formal verification. Conversely, the fact that the ingredients are classical does not remove the need to prove the comparison maps and to retain the right hypotheses.

\clearpage\begingroup\setlength{\parskip}{0pt}\tableofcontents\endgroup
\clearpage\pagenumbering{arabic}
\input{chapters/01_views}
\input{chapters/02_obstructions}
\partnote{Resolution, repair, and retained detail}{Resolution gives a concrete family of views. Its compatibility differential connects exact maps, protected information, and residual coordinates. Recovering every finite coordinate is then separated from convergence in a chosen norm.}
\section{Resolution towers and exact transport}\label{sec:towers}
\subsection{The sequential model}
A tower in a category \(\mathcal C\) is a sequence
\[
 \cdots\xrightarrow{\rho_2}X_2\xrightarrow{\rho_1}X_1
 \xrightarrow{\rho_0}X_0.
\]
For sets, a point is a sequence \(x_n\in X_n\) with \(\rho_n(x_{n+1})=x_n\). Composite restrictions turn the tower into a view system on \(\N\), with an arrow \(i\to j\) when \(i\le j\). Its section space is the same coherent point type, with the same coordinates.

A candidate transport from \(X\) to \(Y\) consists of maps \(f_n:X_n\to Y_n\). It is exact when
\begin{equation}\label{eq:exacttower}
 \rho^Y_n f_{n+1}=f_n\rho^X_n\quad\text{for all }n.
\end{equation}

\begin{proposition}[The exact tower category]\label{KS16}
Identity maps and composites of exact transports are exact. An exact transport induces a map of coherent point spaces by \(x\mapsto(f_nx_n)_n\), preserving identities and composition.
\end{proposition}
\begin{proof}
For two exact transports, substitute \(\rho^Z_ng_{n+1}=g_n\rho^Y_n\) and \(\rho^Y_nf_{n+1}=f_n\rho^X_n\) into the composite square. For a coherent point,
\[
 \rho^Y_n f_{n+1}x_{n+1}
 =f_n\rho^X_nx_{n+1}=f_nx_n.
\]
Identity, associativity, and composition on points are checked coordinatewise.
\end{proof}
\source{KS16}
This is a category on a fixed schedule. A full procategory allows additional reindexing and equivalence relations on maps. Those constructions must be supplied rather than identified with the fixed schedule by terminology.

\subsection{The actual residual of a transport}
Assume \(\mathcal C\) is preadditive, so differences of maps make sense. The tower differential is
\[
 (\partial f)_n=\rho^Y_nf_{n+1}-f_n\rho^X_n.
\]
Its target has components \(\Hom(X_{n+1},Y_n)\), not \(\Hom(X_n,Y_n)\).

\begin{proposition}[The transport Leibniz identity]\label{KS17}
For composable candidates,
\[
 \partial(gf)_n=(\partial g)_nf_{n+1}+g_n(\partial f)_n.
\]
In particular exactness is equivalent to vanishing of the differential.
\end{proposition}
\begin{proof}
Expand the right side. The terms \(-g_n\rho^Y_nf_{n+1}\) and \(+g_n\rho^Y_nf_{n+1}\) cancel, leaving \(\rho^Z_ng_{n+1}f_{n+1}-g_nf_n\rho^X_n\). This is the differential of the composite. The exactness statement is \eqref{eq:exacttower}.
\end{proof}
\source{KS17}
There is no additional mixed defect term in this identity. Such a term would require a different definition of the defect or a different expansion. Typed cancellation is a useful safeguard against importing an attractive but incorrect error formula.

\subsection{Cofinal views preserve points, not every measurement}
Let \(\phi:\N\to\N\) be monotone and unbounded. The sampled tower uses levels \(X_{\phi(n)}\) and their composite restrictions.

\begin{theorem}[Cofinal reconstruction of points]\label{KS18}
Sampling induces an equivalence between the coherent points of \(X\) and of its sampled tower.
\end{theorem}
\begin{proof}
For a coherent sampled point \(y\) and an arbitrary level \(k\), choose \(n\) with \(k\le\phi(n)\) and restrict \(y_n\) to \(X_k\). Two choices agree after passing to a later sampled stage, using monotonicity and coherence. These recovered coordinates satisfy every bonding equation. Restricting a full point and recovering it gives its original coordinates, while sampling the recovered family gives \(y\).
\end{proof}
\source{KS18}
This theorem is about exact coherent points. Raw residual size is a different quantity. On identity scalar towers, \(f_n=(-1)^n\) has nonzero fine differences, but its even subsequence is constant. The coarse residual vanishes because adjacent fine defects cancel. A cofinal equivalence of limits does not imply equality of residual norms or an absence of aliasing in uncorrected data.

\plate{03_tower}{Exact transport commutes with every adjacent restriction. A cofinal subsequence retains coherent points, but a raw defect can cancel during coarsening. The retained coarse value alone need not record the omitted fine discrepancy.}{fig:tower}

\chap{Banach repair and constructed split coordinates}\label{sec:split}
\subsection{An actual complete candidate space}
Let \(X_n,Y_n\) be Banach spaces over a complete nontrivially normed field, with continuous linear bonding maps and uniform bounds \(\norm{\rho^X_n}\le B_X\), \(\norm{\rho^Y_n}\le B_Y\). Use the spaces of bounded operator families
\[
 E=\ell^\infty_n\mathcal L(X_n,Y_n),\qquad
 R=\ell^\infty_n\mathcal L(X_{n+1},Y_n).
\]
The source formalization derives the required completeness from the operator spaces and their bounded product. The same differential satisfies
\[
 \norm{\partial f}_\infty\le(B_X+B_Y)\norm f_\infty.
\]
Thus it is a continuous linear map, and evaluation at a fixed level is continuous. These facts connect algebraic compatibility to an analytic limit argument.

\subsection{From a correction operator to a projection}
The next statement is general. Let \(E\) be Banach, let \(d:E\to R\), \(p:E\to S\), and \(A:R\to E\) be continuous linear maps. Suppose \(pA=0\) and, for \(0\le q<1\),
\begin{equation}\label{eq:contraction}
 \norm{df-dA(df)}\le q\norm{df}\qquad(f\in E).
\end{equation}
The hypothesis concerns residuals of actual candidates. It need not hold for every independent vector in \(R\).

\begin{theorem}[Constructed geometric repair]\label{KS19}
Starting from \(f^{(0)}=f\), define
\[
 f^{(m+1)}=f^{(m)}-A(df^{(m)}).
\]
There is a bounded linear projection \(P_A:E\to\ker d\) with
\[
 f^{(m)}\longrightarrow P_Af,\qquad p(P_Af)=p(f),
\]
and
\[
 \norm{f^{(m)}-P_Af}
 \le\frac{\norm A}{1-q}\norm{df^{(m)}}.
\]
It fixes every exact candidate and satisfies
\[
 \norm{P_A}\le1+\frac{\norm A\norm d}{1-q}.
\]
\end{theorem}
\begin{proof}
Equation~\eqref{eq:contraction} gives geometric decay of the residuals. The step size is at most \(\norm A\norm{df^{(m)}}\), so the series of movements converges and the sequence is Cauchy. Completeness supplies a limit. Continuity of \(d\) gives zero residual, and \(pA=0\) makes every iterate preserve \(p(f)\). Summing the tail of the geometric movement estimate gives the displayed error bound.

Each iterate is linear in its initial value. Addition and scalar multiplication commute with the limit, so the limit map is linear. The estimate at \(m=0\), followed by the triangle inequality, gives its operator bound. If \(df=0\), the iteration is stationary. Therefore its image is exactly \(\ker d\), it fixes that image, and it is idempotent.
\end{proof}
\source{KS19}
This theorem constructs the dynamics from \(A\). It does not construct a suitable \(A\) for every residual problem. That missing input can, however, be derived for a useful class of actual towers.

\subsection{A right inverse from target splittings}
Suppose each target bonding map has a continuous linear section
\[
 s_n:Y_n\to Y_{n+1},\qquad \rho^Y_ns_n=I,\qquad \norm{s_n}\le S.
\]
Assume the propagation gap
\begin{equation}\label{eq:splitgap}
 SB_X<1,\qquad K=\frac{S}{1-SB_X}.
\end{equation}
For \(r\in R\), recursively set
\begin{equation}\label{eq:greenrecurrence}
 (Ar)_0=0,\qquad
 (Ar)_{n+1}=s_n\bigl(r_n+(Ar)_n\rho^X_n\bigr).
\end{equation}

\begin{theorem}[Split tower correction]\label{KS20}
The recurrence defines a continuous linear map \(A:R\to E\) with
\[
 \partial A=I_R,\qquad \norm A\le K.
\]
The operator \(P=I-A\partial\) is an exact projection and
\[
 (Pf)_0=f_0,\qquad
 \norm{f-Pf}_\infty\le K\norm{\partial f}_\infty.
\]
\end{theorem}
\begin{proof}
Inductively, if \(\norm{(Ar)_n}\le K\norm r_\infty\), then
\[
 \norm{(Ar)_{n+1}}
 \le S(1+B_XK)\norm r_\infty=K\norm r_\infty.
\]
The base case is zero. The recurrence is linear, so this uniform bound bundles it into a bounded linear operator. Applying \(\rho^Y_n\) to \eqref{eq:greenrecurrence} and using the section identity gives
\[
 (\partial Ar)_n=r_n+(Ar)_n\rho^X_n-(Ar)_n\rho^X_n=r_n.
\]
Consequently \(\partial P=0\), \(P^2=P\), and the error is \(A\partial f\). The initial correction vanishes.
\end{proof}
\source{KS20}
The surjectivity is onto all bounded residual families. It is stronger than the assertion that some exact output can be assigned to each candidate. The gap and the chosen splittings are genuine inputs.

\subsection{Protected information beyond the seed}
When the target grows, its new kernel directions are independent information. Define
\[
 N_n(f)=(I-s_n\rho^Y_n)f_{n+1},\qquad
 (p_Lf)_0=f_0,\quad(p_Lf)_{n+1}=N_n(f).
\]
Each \(N_n(f)\) lands in \(\ker\rho^Y_n\). The observable \(p_L\) retains both the seed and every new output component, relative to the chosen sections.

\begin{theorem}[Split coordinate theorem]\label{KS21}
The observable is a bounded projection with \(p_LA=0\) and \(p_LP=p_L\). It separates exact candidates. If
\[
 \mathcal S_L=\{s\in E:p_Ls=s\},
\]
then
\[
 \Phi:E\simeq_L\mathcal S_L\times R,
 \qquad \Phi(f)=(p_Lf,\partial f),
 \qquad \Phi^{-1}(s,r)=Ps+Ar.
\]
Repair becomes \((s,r)\mapsto(s,0)\). It is the unique exact repair preserving \(p_Lf\).
\end{theorem}
\begin{proof}
The operator \(I-s_n\rho^Y_n\) is a projection onto the target kernel. This proves idempotence of \(p_L\), and its norm is at most \(1+SB_Y\). Every correction in \eqref{eq:greenrecurrence} at a positive level lies in the range of the corresponding section, so its innovation is zero. Hence \(p_LA=0\) and \(p_LP=p_L\).

If \(f\) is exact, then
\[
 f_{n+1}=s_nf_n\rho^X_n+N_n(f).
\]
Its seed and innovations determine the family by induction. Therefore two exact families with the same observable agree. More generally, if \(f,g\) have the same observable and residual, apply this fact to the exact difference \(f-g\).

For \((s,r)\in\mathcal S_L\times R\), the candidate \(Ps+Ar\) has observable \(s\) and residual \(r\). Conversely \(f\) and \(P(p_Lf)+A\partial f\) have the same two coordinates, so they agree. Boundedness of \(p_L\), \(\partial\), \(P\), and \(A\) makes both inverse maps continuous. Finally, zeroing the residual retains \(s\), and separation proves uniqueness.
\end{proof}
\source{KS21}
This is an actual analytic coordinate equivalence, not generally an isometry. The protected coordinate depends on the declared splitting. It should not be called an intrinsic semantic coordinate without a further comparison theorem.

\plate{04_split}{In the split coordinate chart, a candidate separates into retained information and compatibility residual. The exact subspace is the zero residual slice. The chart and its protected coordinate depend on the chosen target sections.}{fig:split}

\subsection{A growing example and the gap boundary}
Use \(X_n=Y_n=\R^{n+1}\) with maximum norms. Let \(\pi_n\) delete the last coordinate, take \(\rho^X_n=\tfrac12\pi_n\), \(\rho^Y_n=\pi_n\), and let \(s_n\) append zero. Then \(S=B_Y=1\), \(B_X=1/2\), and \(K=2\).

\begin{example}[Repair of the identity family]\label{KS22}
The bounded candidate \(f_n=I\) is not exact. Its repair is
\[
 (Pf)_n=\operatorname{diag}(2^{-n},2^{-(n-1)},\ldots,2^{-1},1).
\]
The new last output row is retained at every stage while inherited rows are propagated by the factor \(1/2\).
\end{example}
\begin{proof}
At level zero the value is one. Exactness forces every inherited row at the next level to be one half of its previous value. Preservation of the innovation forces the new row to remain the corresponding identity row. Induction gives the displayed diagonal.
\end{proof}
\source{KS22}
For constant identity towers, by contrast, \((\partial f)_n=f_{n+1}-f_n\). The constant residual \(r_n=1\) has no bounded primitive, since any primitive satisfies \(f_n=f_0+n\). Here \(SB_X=1\). This counterexample excludes discarding the gap indiscriminately. It does not prove that the gap is necessary for every individual tower. Exact repair preserving the initial coordinate still exists in this constant example, by replacing a candidate with its initial constant family. Surjectivity onto all bounded residuals does not follow from that repair.
\source{KS23}

\subsection{What is canonical}
For fixed splitting and protected observable, the repair above is unique. Changing the splitting can change both the observable and the repaired answer. In the growing example, replace \(s_n(y)=(y,0)\) by \(s_n'(y)=(y,y_0)\). At the first positive level, repair of the identity becomes
\[
 \begin{pmatrix}1/2&0\\-1/2&1\end{pmatrix}
 \quad\text{instead of}\quad
 \begin{pmatrix}1/2&0\\0&1\end{pmatrix}.
\]
Both are exact continuations for their respective innovation profiles and have the same seed. Thus a unique answer for a declared preservation problem is not an answer independent of all preservation choices.
\source{KS24}

\chap{Multiscale residuals and the limits of reconstruction}\label{sec:cascade}
\subsection{A split coarsener retains its missing detail}
Let \(C:R\to\overline R\) and \(J:\overline R\to R\) be linear maps with \(CJ=I\). Put \(H=I-JC\). Then
\[
 r=JCr+Hr,\qquad CH=0,\qquad H^2=H.
\]
The two coordinates \((Cr,Hr)\) retain all information. The coarse coordinate alone generally does not.

The actual dyadic tower coarsener has the typed formula
\begin{equation}\label{eq:dyadic}
 (Cr)_n=\rho^Y_{2n}r_{2n+1}+r_{2n}\rho^X_{2n+1}.
\end{equation}
It is obtained by telescoping two adjacent compatibility squares. With target sections, a lift is
\[
 (Jc)_{2n}=0,\qquad (Jc)_{2n+1}=s_{2n}c_n.
\]
The section identity gives \(CJ=I\), so the missing detail is explicitly \(Hr\), not an unspecified reconstruction error. The original normed construction imposes the bounds needed to make these maps bounded on its operator family spaces.
\source{KS25}

\subsection{A whole cascade}
At scale \(j\), let \(C_j:R_j\to R_{j+1}\) and \(J_j:R_{j+1}\to R_j\) be bounded linear maps with \(C_jJ_j=I\). For \(x\in R_0\), define
\[
 x_0=x,\quad x_{j+1}=C_jx_j,\quad
 L_0=I,\quad L_N=J_0\cdots J_{N-1},
\]
\[
 h_j=x_j-J_jx_{j+1},\qquad Q_Nx=L_Nx_N,
 \qquad S_N=I-Q_N.
\]
The map \(Q_N\) restores the surviving coarse remainder to the original space. It is this restored remainder, not just \(x_N\), that controls reconstruction.

\begin{theorem}[Finite reconstruction and nested remainders]\label{KS26}
For every \(N\),
\[
 x=Q_Nx+\sum_{j<N}L_jh_j,
 \qquad S_Nx=\sum_{j<N}L_jh_j.
\]
Each \(Q_N\) is a projection, and for \(N\le M\),
\[
 Q_NQ_M=Q_MQ_N=Q_M.
\]
The vectors with all details zero are exactly the persistent subspace
\[
 \mathcal P_\infty=\{x:Q_Nx=x\text{ for every }N\}.
\]
\end{theorem}
\begin{proof}
The one scale identity gives \(L_jh_j=Q_jx-Q_{j+1}x\). Summing telescopes from \(Q_0=I\). A coarsening composite followed by its matching restoration cancels because \(C_jJ_j=I\). This proves idempotence and both nested projection identities. If all details vanish, successive remainders are equal and hence all equal \(x\). The converse follows from their successive differences.
\end{proof}
\source{KS26}

\plate{05_cascade}{A split cascade retains every discarded component as a detail. The finite sum is always exact after the restored coarse remainder is included. An infinite detail expansion requires that remainder to vanish in the topology being used.}{fig:cascade}

\subsection{Convergence has its own theorem}
Call \(x\) norm recoverable when \(Q_Nx\to0\). This is equivalent to \(S_Nx\to x\), but it does not follow merely from vanishing persistent subspace. Let
\[
 \mathcal F=\bigcup_N\ker Q_N
\]
be the finite core.

\begin{theorem}[Density criterion for norm recovery]\label{KS27}
Suppose \(\sup_N\norm{Q_N}\le K<\infty\). Then the norm recoverable vectors are exactly \(\overline{\mathcal F}\). In particular all vectors are norm recoverable if and only if the finite core is dense.
\end{theorem}
\begin{proof}
If \(Q_Nx\to0\), then \(S_Nx\to x\), and \(Q_NS_Nx=0\), so each approximant lies in the finite core. Conversely choose \(y\in\ker Q_M\) close to \(x\). For \(N\ge M\), nesting gives \(Q_Ny=0\), so
\[
 \norm{Q_Nx}\le K\norm{x-y}.
\]
Choose the approximation first and then the depth. This proves convergence for every \(x\) in the closure.
\end{proof}
\source{KS27}
There are other sufficient conditions. If the transported details are absolutely summable in the Banach space \(R_0\), their sum exists and the remaining boundary lies in \(\mathcal P_\infty\). If that persistent subspace is zero, the sum equals \(x\). Alternatively,
\[
 \norm{x-S_Nx}\le\left(\prod_{j<N}\norm{J_j}\right)\norm{x_N}
\]
gives a direct sufficient decay test. A decreasing coarse norm by itself is not enough when restoration gains grow.

\subsection{The actual dyadic result}
For the source's dyadic residual hierarchy, the persistent residual is zero. Zero detail at a scale reveals every even coordinate, while an odd coordinate is passed to a smaller index at the next scale. Strong induction on the coordinate index therefore proves that a residual with all details zero is zero.

\begin{theorem}[Dyadic determination and finite coordinate recovery]\label{KS28}
In the actual dyadic system, all detail coordinates jointly determine the residual. Moreover,
\[
 k+1<2^N\quad\Longrightarrow\quad
 (Q_Nr)_k=0,\qquad(S_Nr)_k=r_k.
\]
Thus every fixed residual coordinate is recovered exactly at a sufficiently large finite depth. The statement does not assert convergence in the supremum norm over all coordinates.
\end{theorem}
\begin{proof}[Proof outline]
At one scale, the hidden even coordinate is the original even coordinate. An odd hidden equation expresses the remaining coordinate through the next scale. Repeated descent reduces an index \(k\) to \(\lfloor k/2\rfloor\). After \(N\) steps with \(k+1<2^N\), no unrecovered coordinate remains at that index. This proves the displayed support statement, which is formalized with the dependent tower operator types. Applying it to a residual with zero detail at every scale forces every coordinate to vanish.
\end{proof}
\source{KS28}
The strict inequality matters. At the final coordinate of a dyadic block, the remainder can still be nonzero. This is a support theorem, not an argument based on finitely many numerical samples.

\subsection{Two boundary examples}
On \(\ell^\infty(\N,\R)\), let \(C\) remove the first coordinate and let \(J\) prepend zero. Then \(Q_Nx\) is the tail of \(x\), placed at its original coordinates with zeros before \(N\). Every fixed coordinate is eventually zero in the remainder, and \(\norm{Q_N}\le1\). Nevertheless the constant sequence \(\mathbf1\) satisfies \(\norm{Q_N\mathbf1}=1\) for every \(N\).

\begin{proposition}[The recovery domain is \(c_0\)]\label{KS29}
For this shift cascade,
\[
 S_Nx\to x\text{ in }\ell^\infty
 \quad\Longleftrightarrow\quad x_n\to0.
\]
For \(x_n=q^n\), \(0\le q<1\), the error is exactly \(q^N\).
\end{proposition}
\begin{proof}
The error norm is \(\sup_{n\ge N}|x_n|\). It tends to zero precisely when the sequence tends to zero. For a geometric sequence its largest tail value is the first one, \(q^N\).
\end{proof}
\source{KS29}
As a second example, take scalar spaces with \(C=\tfrac12I\), \(J=2I\). The coarse value \(x_N=2^{-N}x\) tends to zero, but restoration gives \(Q_Nx=x\). Every detail is zero and the entire input is persistent. The example belongs to this scalar cascade, not to the actual dyadic hierarchy, whose persistent subspace was proved to vanish.

\begin{insight}
Packets of retained detail are useful because they make omission explicit. They do not automatically make an infinite reconstruction convergent. Coordinate separation, bounded reconstruction, density, summability, and attainment remain different mathematical properties.
\end{insight}

\partnote{Symmetry and effective descent}{A set of solutions forgets the witnesses relating equivalent solutions. The next layer retains those witnesses, their composition, and their automorphisms. Gluing then has to respect the overlap identifications themselves.}
\section{Nonabelian coherence and relative trivialization}\label{sec:gauge}
\subsection{Coherent transition labels}
Fix groups \(G_v\) and restriction homomorphisms \(f^*:G_w\to G_v\), with the same identity and composition laws as before. The groups need not commute and the homomorphisms need not be bijective.

\begin{definition}[Cocycle]
A cocycle is a family \(c_f\in G_v\) for every \(f:v\to w\), satisfying
\[
 c_{1_v}=1,\qquad c_{g\circ f}=c_f f^*(c_g).
\]
It determines a twisted view system with maps
\[
 T_c(f)(x)=c_f f^*(x).
\]
A section of the twisted system satisfies \(x_v=c_f f^*(x_w)\).
\end{definition}
The cocycle law says that the transition rules compose coherently. It does not assert that the twisted system has a section. The coherent transition system and a global trivialization of that system are different objects.

For a family \(a=(a_v)\in\prod_vG_v\), define
\begin{equation}\label{eq:gauge}
 (a\cdot c)_f=a_vc_f[f^*(a_w)]^{-1}.
\end{equation}

\begin{proposition}[Gauge action and its groupoid]\label{KS30}
Equation~\eqref{eq:gauge} defines an action on cocycles. A gauge carrying \(c\) to \(d\) induces an equivalence of twisted section sets by \(x_v\mapsto a_vx_v\). The cocycles and actual gauges form a groupoid: an arrow \(a:c\to d\) followed by \(b:d\to e\) has label \(ba\).
\end{proposition}
\begin{proof}
For composable \(f:v\to w\) and \(g:w\to z\), the middle factors cancel in
\[
 (a\cdot c)_f f^*((a\cdot c)_g)
 =a_vc_f f^*(c_g)[f^*g^*(a_z)]^{-1}.
\]
This is the transformed label at \(g\circ f\). The identity law is similar. Direct substitution gives \(a\cdot(b\cdot c)=(ab)\cdot c\), with identity gauge acting trivially. If \(x\) is twisted coherent, then
\[
 (a\cdot c)_f f^*(a_wx_w)=a_vc_f f^*(x_w)=a_vx_v.
\]
The inverse gauge gives the inverse section map. Identity arrows, composition, and inverses in the groupoid now follow from the group action laws.
\end{proof}
\source{KS30}
The action groupoid \(Z^1(\V;G)//\prod_vG_v\) retains comparison witnesses. Its orbit set \(Z^1(\V;G)/\prod_vG_v\) retains only isomorphism classes. There is no reason for the orbit set to be a group in this nonabelian setting.

\subsection{Automorphisms and ambiguity are not the same question}
The stabilizer of a cocycle is the group of families with
\[
 a_v=c_f f^*(a_w)c_f^{-1}.
\]
If one gauge \(a:c\to d\) is known, all gauges from \(c\) to \(d\) are uniquely of the form \(ah\), with \(h\in\Stab(c)\). Indeed \(a^{-1}b\) stabilizes \(c\) for every other gauge \(b\). This is the ordinary Hom torsor principle in a groupoid.

That symmetry group must not be confused with the set of global trivializations of one fixed cocycle. The latter will be parametrized by ordinary coherent group elements, after one trivialization has been selected.

\subsection{The relative obstruction without commutativity}
Prescribe a coherent partial section \(s\) of \(T_c\) on \(U\). The based gauge group is
\[
 \mathcal G_U=\{a\in\prod_vG_v:a_u=1\text{ for }u\in U\}.
\]
Choose a raw family \(q\) equal to \(s\) on \(U\). Identity elements can be used elsewhere. Normalize the labels by
\[
 c^q_f=q_v^{-1}c_f f^*(q_w).
\]
Its labels on arrows internal to \(U\) are one. Define the pointed orbit obstruction
\[
 \mathfrak o_U(c,s)=[c^q]\in Z^1(\V;G)/\mathcal G_U,
\]
whose neutral point is the class of the trivial cocycle.

\begin{theorem}[Nonabelian relative extension]\label{KS31}
The obstruction is independent of the raw lift. It is neutral exactly when \(s\) extends to a global section of \(T_c\).
\end{theorem}
\begin{proof}
If \(t\) is another raw lift, the family \(t^{-1}q\) is based and carries \(c^q\) to \(c^t\). Thus the orbit does not depend on the lift.

If a based gauge \(b\) carries \(c^q\) to the trivial cocycle, set \(x_v=q_vb_v^{-1}\). The gauge equation becomes \(x_v=c_f f^*(x_w)\), and \(b|_U=1\) gives \(x|_U=s\). Conversely, an extension \(x\) gives the based gauge \(b=x^{-1}q\), which trivializes \(c^q\). All products retain the displayed order. No commutativity has been used.
\end{proof}
\source{KS31}
The orbit is defined from actual cocycles and an actual group action. It is not a Boolean test defined to be neutral precisely when an extension exists. Deciding orbit membership in arbitrary groups remains a separate problem.

Let
\[
 H_U=\{h:f^*(h_w)=h_v\text{ for every }f,\ h|_U=1\}.
\]

\begin{theorem}[The nonabelian extension torsor]\label{KS32}
Given one extension \(x_0\), the map \(h\mapsto x_0h\) is an equivalence from \(H_U\) to the extension set. The extension is unique if and only if the obstruction is neutral and \(H_U\) is trivial.
\end{theorem}
\begin{proof}
If \(h\) is coherent, then
\[
 c_f f^*((x_0)_wh_w)=(x_0)_vh_v,
\]
so right multiplication gives an extension. If \(x\) is another extension, then
\[
 (x_0)_v^{-1}x_v
 =[c_f f^*((x_0)_w)]^{-1}[c_f f^*(x_w)]
 =f^*((x_0)_w^{-1}x_w).
\]
Thus \(x_0^{-1}x\in H_U\) and is the inverse parameter. The uniqueness statement follows as in the additive case.
\end{proof}
\source{KS32}
The additive theory is recovered by applying a multiplicative type tag to an additive view system and choosing the trivial cocycle. The normalized label then becomes \(f^*q_w-q_v\). The source proves agreement of the actual extension sets and vanishing criteria. It does not identify the whole nonabelian orbit set with every element of the larger additive arrow residual quotient.

\subsection{A small noncommutative model}
\begin{example}[A \(C_2\) loop with \(S_3\) coefficients]\label{KS33}
Take one view with a loop \(\sigma\) satisfying \(\sigma^2=1\), and constant coefficients in a group \(H\). A cocycle is determined by an involution \(p\in H\). Gauge equivalence is conjugacy, and its automorphisms form the centralizer of \(p\).

For \(H=S_3\), the involutions are the identity and the three transpositions. There are exactly two orbit classes. The identity cocycle has automorphism group \(S_3\), which is noncommutative. A nonidentity involution gives a coherent cocycle whose regular twisted section equation \(x=px\) has no solution.
\end{example}
\begin{proof}
The only nontrivial cocycle relation is \(p^2=1\). At a single view the gauge action is \(p\mapsto apa^{-1}\). All three transpositions are conjugate and none is conjugate to the identity. The stabilizer is therefore the centralizer. Finally \(x=px\) implies \(p=1\) after multiplication by \(x^{-1}\).
\end{proof}
\source{KS33}
The example makes an important boundary visible. A cocycle can exist and glue as an object even when it admits no global trivializing section. Stackhood will concern the former gluing problem, not a claim that every cocycle becomes trivial.

\plate{06_groupoid}{The orbit set sees two classes in the \(C_2\)-to-\(S_3\) example. The groupoid also retains the arrows between conjugate transpositions and all automorphisms. Collapsing the arrows loses information that matters in descent.}{fig:groupoid}

\chap{Effective gluing of objects and their identifications}\label{sec:descent}
\subsection{Why pairwise agreement is not the right datum}
For objects with symmetry, local descriptions need not be literally equal on overlaps. They may instead come with specified isomorphisms. Those isomorphisms must satisfy their own composition law. It is not enough to know that some isomorphism exists for each pair.

Let \(U_i\) be full patches in \(\V\). Assume the composition cover condition
\begin{equation}\label{eq:compositioncover}
 \forall(v\xrightarrow f w\xrightarrow g z),\quad
 \exists i:\ v,w,z\in U_i.
\end{equation}
By inserting identity arrows, this implies coverage of objects and arrows.

A descent datum consists of local cocycles \(c_i\) and overlap gauges \(a_{ij}(v)\) with
\begin{align}
 c_j(f)&=a_{ij}(v)c_i(f)[f^*a_{ij}(w)]^{-1},\label{eq:localgauge}\\
 a_{ik}(v)&=a_{jk}(v)a_{ij}(v).\label{eq:triple}
\end{align}
The triple equation implies \(a_{ii}=1\) and \(a_{ji}=a_{ij}^{-1}\). A morphism of descent data is a family of local gauges \(r_i\) satisfying
\[
 r_j(v)a_{ij}(v)=a'_{ij}(v)r_i(v).
\]
These objects and arrows form a groupoid.

\begin{theorem}[Effective cocycle descent]\label{KS34}
Under \eqref{eq:compositioncover}, restriction is an equivalence between the global cocycle gauge groupoid and the groupoid of descent data. The overlap gauges supplied in the datum are realized, not discarded.
\end{theorem}
\begin{proof}
Choose a reference patch \(\kappa(v)\) through each view and set
\[
 b_i(v)=a_{\kappa(v),i}(v).
\]
Equation~\eqref{eq:triple} gives \(b_j=a_{ij}b_i\). Normalize the local labels by
\[
 \widetilde c_i(f)=b_i(v)^{-1}c_i(f)f^*(b_i(w)).
\]
Substituting \eqref{eq:localgauge} and \(b_j=a_{ij}b_i\) shows that these normalized labels agree literally on overlaps. They can therefore be assembled into labels \(c_f\). For every composable pair, a patch containing all three endpoints verifies the cocycle law for the assembled labels. Thus \(c\) is a global cocycle.

Undoing the normalization gives
\[
 b_i(v)c_f[f^*(b_i(w))]^{-1}=c_i(f),\qquad b_j=a_{ij}b_i.
\]
So the local identifications implement the original overlap gauges. This constructs an object whose restriction is isomorphic to the specified descent datum.

For morphisms between restrictions of two global cocycles, the overlap gauges are identities. The morphism condition reduces to equality of the local gauges on overlaps. They assemble uniquely into a global gauge, and arrow coverage verifies its compatibility with the cocycles. Restriction and assembly are inverse on Hom sets. The restriction functor is therefore fully faithful and essentially surjective, hence an equivalence of groupoids.
\end{proof}
\source{KS34}

\subsection{The correct uniqueness statement}
Suppose \((c,b_i)\) and \((c',b'_i)\) realize the same descent datum. The only possible compatible comparison is
\[
 g_v=b'_i(v)^{-1}b_i(v),
\]
using any patch through \(v\). The overlap identities make it independent of the patch, and \(b'_ig=b_i\). It is therefore the unique global comparison respecting the fixed local identifications.

This is not a statement that the underlying cocycle has no automorphisms. Those automorphisms can reappear when the local identifications are forgotten. Nor does the construction choose one representative independently of every reference patch. It proves that the choices are related by the appropriate unique comparison.

\subsection{A failed triple}
Let \(t\ne1\) be an involution and ask for frames \(b_0,b_1,b_2\) satisfying
\[
 b_1=tb_0,\qquad b_2=tb_1,\qquad b_2=tb_0.
\]
The first two give \(b_2=b_0\), while the third forces \(t=1\). Each displayed transition is individually invertible, but they cannot be the overlap identifications of a global frame. The defect lies in the triple relation.

\plate{07_cover}{Object coverage, arrow coverage, and composition coverage are distinct. The proper two-view patches on a three-view chain see every arrow but no patch sees the relation \(0\to1\to2\). The four proper three-view patches on a four-view chain do cover every composable pair.}{fig:cover}

\chap{A genuine stack on the composition site}\label{sec:stack}
The preceding theorem can be placed in a standard site, rather than called stack-like. We use the usual effective-descent meaning of a stack~\cite{stacks-stack,mathlib-stack}. The base of this site is the inclusion poset of full patches
\[
 \mathcal P(\V)=(\operatorname{Set}(\Ob\V),\subseteq).
\]
This patch category is different from the original category of views.

\subsection{The topology is defined without assuming effectivity}
For a composable pair \(q=(v\xrightarrow f w\xrightarrow g z)\), let its support be \(|q|=\{v,w,z\}\). A sieve \(\mathcal R\) on \(W\) is covering when
\begin{equation}\label{eq:site}
 |q|\subseteq W\quad\Longrightarrow\quad
 (|q|\hookrightarrow W)\in\mathcal R
 \quad\text{for every composable }q.
\end{equation}
This condition depends only on the view category, not on the coefficient groups or on a promised global realization.

\begin{proposition}[The composition topology]\label{KS35}
Condition~\eqref{eq:site} defines a Grothendieck topology \(J_{\mathrm{comp}}\). The sieve generated by a family of inclusions \(U_i\subseteq W\) is covering exactly when the family covers every composition support in \(W\).
\end{proposition}
\begin{proof}
The maximal sieve contains all support inclusions. If \(Z\subseteq W\), every support in \(Z\) is also a support in \(W\), so its inclusion belongs to the pullback of a covering sieve. For transitivity, let \(\mathcal R\) cover \(W\), and let another sieve \(\mathcal T\) cover locally over every member of \(\mathcal R\). The support inclusion \(|q|\hookrightarrow W\) belongs to \(\mathcal R\). Applying the local covering condition on \(|q|\) to the same composable pair puts the identity inclusion of \(|q|\) into the pulled back \(\mathcal T\). Thus \(|q|\hookrightarrow W\) belongs to \(\mathcal T\).

In the poset of subsets, an inclusion of a support factors through a member of the family precisely when that member contains the support. Membership in the generated sieve is exactly this factorization. This proves the final characterization.
\end{proof}
\source{KS35}
Intersections of a covering family with a subpatch remain covering. Replacing each member by a covering family again gives a cover. Identities imply coverage of single views and arrows. The empty patch admits the empty cover.

\subsection{The actual patch pseudofunctor}
On a patch \(W\), use
\[
 \mathfrak K_G(W)=Z^1(\V|_W;G|_W)//\prod_{v\in W}G_v.
\]
Restriction keeps the existing cocycle values and gauge values on the smaller patch. It first gives a strict category valued functor on the opposite patch poset. Its standard pseudofunctor form has identity comparison isomorphisms, with the required coherence laws.

\begin{theorem}[Stackhood on the composition site]\label{KS36}
The patch pseudofunctor \(\mathfrak K_G\) is a stack for \(J_{\mathrm{comp}}\). It satisfies the standard condition of descent for morphisms and effective descent for objects on every covering sieve.
\end{theorem}
\begin{proof}
For a covering family, Theorem~\ref{KS34} gives a concrete equivalence of groupoids. The standard descent interface allows comparison morphisms on arbitrary common subpatches, not just pointwise labels. Evaluate these comparison morphisms on singleton patches. Their naturality under restriction identifies those values with their values on every larger common subpatch. The standard cocycle law becomes \eqref{eq:triple}. Conversely, the concrete effectuation gives the required standard local identifications, and the same naturality proves every common subpatch equation.

The morphism construction retains all Hom sets and their inverse gluing maps. Thus the standard restriction functor is fully faithful and essentially surjective. Finally, index an arbitrary covering sieve by its own arrows. Every composition support is one of those arrows by definition, so that family satisfies \eqref{eq:compositioncover}. The family equivalence proves effective descent for every covering sieve, which is the standard stack condition.
\end{proof}
\source{KS36}
The principal source declaration is \code{GroupViewSystem.patch_isStack}. Its target is mathlib's \code{Pseudofunctor.IsStack}, not a project predicate that repeats the conclusion as an assumption.

\subsection{The cover is nontrivial, but specific}
On a four-view chain, take the four patches obtained by deleting one vertex. All patches are proper, yet every composable pair has at most three endpoints and lies in one of them. The theorem genuinely glues pieces smaller than the whole object.

On the three-view chain, the analogous proper patches cover each object and each arrow but miss the support of \(0\to1\to2\). They do not form a cover in this topology. Consequently the theorem is neither vacuous nor an unrestricted gluing principle.

\begin{scopebox}
\textbf{Exact scope.} This is an ordinary stack of strict group valued cocycle realizations on a specified combinatorial site. It is not an infinity-stack, a universal heterogeneous atlas, or an arithmetic existence theorem. The stack theorem glues coherent cocycles with their overlap isomorphisms. It does not assert that every resulting cocycle has a global trivialization.
\end{scopebox}

\chap{Equivalent presentations and identified boundaries}\label{sec:relative}
\subsection{Transporting the same question}
Let \(H:\V\to\W\) be part of a category equivalence, with inverse \(L\) and counit \(\varepsilon:HL\Rightarrow\id_\W\). Pull a view system \(F\) on \(\W\) back along \(H\). Sampling a global section gives \((H^*x)_v=x_{Hv}\).

\begin{proposition}[Equivalence of section spaces]\label{KS37}
Pullback induces \(\Gamma(F)\simeq\Gamma(H^*F)\). If \(U\subseteq\Ob\W\) is closed under isomorphism and \(U'=H^{-1}U\), it also induces
\[
 \Ext_U^F(s)\simeq\Ext_{U'}^{H^*F}(H^*s).
\]
Existence and uniqueness of the prescribed extension are preserved and reflected.
\end{proposition}
\begin{proof}
For a section \(y\) of the pulled back system, use \(t_w=\varepsilon_w^{-1}:w\to HLw\) and set
\[
 (Ry)_w=F(t_w)(y_{Lw}).
\]
Naturality and coherence prove that this is a section. The unit and counit identities show that restriction and recovery are inverse.

For a prescribed domain, \(w\in U\) implies \(HLw\in U\) because these objects are isomorphic. Thus the partial equation for \(y\) applies at \(Lw\). Coherence of \(s\) along \(t_w\) gives \((Ry)_w=s_w\). Hence the same inverse maps restrict to the extension fibers. Any equivalence of sets preserves and reflects nonemptiness and the singleton property.
\end{proof}
\source{KS37}
The coefficients really are pulled back. This does not compare unrelated coefficient systems simply because their diagrams look alike. Isomorphism closure is a sufficient general hypothesis for the relative statement. The theorem does not assert that every unsaturated domain fails.

For group coefficients, the invisible coherent groups \(H_U\) are transported by a multiplicative equivalence, and the parameterizations \(x_0h\) commute with pullback. Thus feasibility and ambiguity are both retained. This is more than a comparison of the number of extensions. It is not an assertion that every possible based gauge orbit set is presentation independent.

\subsection{Why a boundary must retain its identification}
Let \(R:\mathcal C\to\mathcal D\) be a functor between groupoids, and fix \(b\in\mathcal D\). A realization over \(b\) should include a chosen isomorphism to \(b\), rather than require literal equality of representatives.

\begin{definition}[Marked realization groupoid]
The objects of \(\mathcal E_b(R)\) are pairs \((x,\phi)\), with \(x\in\mathcal C\) and \(\phi:Rx\xrightarrow{\sim}b\). An arrow \((x,\phi)\to(y,\psi)\) is an arrow \(a:x\to y\) with \(\psi R(a)=\phi\).
\end{definition}
Composition and inverses preserve this equation, so this is a groupoid. It is the explicit groupoid homotopy fiber construction, without invoking a higher categorical interface.

\begin{proposition}[Existence and invisible automorphisms]\label{KS38}
The groupoid \(\mathcal E_b(R)\) is nonempty exactly when \(b\) is in the essential image of \(R\). An automorphism of \((x,\phi)\) is an automorphism \(a\) of \(x\) satisfying \(R(a)=1_{Rx}\).
\end{proposition}
\begin{proof}
The first assertion is the object definition. For the second, the boundary equation is \(\phi R(a)=\phi\). Cancel the isomorphism \(\phi\). Conversely \(R(a)=1\) makes the equation hold.
\end{proof}
\source{KS38}
This retains precisely the symmetries invisible at the identified boundary. It avoids freezing arbitrary representative coordinates and then mistaking the result for an invariant realization problem.

\subsection{Weakly commuting squares are enough}
Suppose \(E:\mathcal C\to\mathcal C'\) and \(K:\mathcal D\to\mathcal D'\) are equivalences, with functors \(R,R'\) and a natural isomorphism
\[
 \sigma:KR\xRightarrow{\sim}R'E.
\]
The square is not required to commute literally.

\begin{theorem}[Invariance of the marked realization problem]\label{KS39}
The map
\[
 (x,\phi)\longmapsto(Ex,K(\phi)\sigma_x^{-1}),\qquad a\longmapsto E(a)
\]
defines an equivalence \(\mathcal E_b(R)\simeq\mathcal E_{Kb}(R')\).
\end{theorem}
\begin{proof}
Naturality of \(\sigma\) verifies the boundary equation for transported arrows. Full faithfulness follows by lifting an arrow through the fully faithful \(E\), then reflecting its boundary equation through the faithful \(K\).

For essential surjectivity, start from \((y,\psi)\). Choose \(x\) and an isomorphism \(i:Ex\to y\). The composite
\[
 \psi\,R'(i)\,\sigma_x:KRx\to Kb
\]
lifts through full \(K\) to an arrow \(\phi:Rx\to b\). Since the target is a groupoid it is an isomorphism. Its transported identification is \(\psi R'(i)\), so \(i\) gives the required marked comparison. The functor is fully faithful and essentially surjective, hence an equivalence.
\end{proof}
\source{KS39}
Apply this to the actual restriction functors between cocycle patch groupoids and the equivalences induced by an equivalent view presentation. The full realization problem, including comparison witnesses and automorphisms, is transported.

Replacing \(b\) by an isomorphic boundary \(b'\) also gives an equivalence, by composing the identification with \(b\to b'\). For \(R=\id\), the marked realization groupoid has center \((b,1_b)\), and between \((x,\phi)\) and \((y,\psi)\) there is exactly one compatible arrow \(\psi^{-1}\phi\). This is a useful example of contractible realization data, not a claim that the original groupoid has no symmetries.

\plate{08_marked}{The chosen boundary identification is part of a realization. A comparison must commute with it. Under an equivalent presentation, the comparison square may commute through an isomorphism rather than a literal equality, and the marked realization groupoid is still transported correctly.}{fig:marked}

\partnote{The geometry of exact realization}{A categorical space of realizations does not come with a numerical distance by default. In a declared linear model, the compatibility equations become a bounded operator. Closed range, holonomy, and conditioning then determine what small residuals actually guarantee.}
\section{Graph views, holonomy, and reconstruction}\label{sec:graph}
\subsection{Actual transports on a finite graph}
Let a finite connected simple graph have vertex set \(V\), root \(o\), and a chosen route \(p_i:o\to i\) for every vertex, with the route at \(o\) empty. Use a common real normed fiber \(H\). An inverse consistent connection assigns continuous linear equivalences \(\tau_{ij}\), with \(\tau_{ji}=\tau_{ij}^{-1}\). Our convention is that \(\tau_{ij}\) carries the value at \(j\) into the coordinates at \(i\).

The candidate and residual spaces are finite products with maximum norms:
\[
 E=\prod_{i\in V}H,\qquad
 R=\prod_{(i,j)\in\vec E}H,\qquad
 (dx)_{ij}=\tau_{ij}x_j-x_i.
\]
Both orientations of an edge are retained. If \(M\) bounds the norms of all measured edge maps, then \(\norm d\le M+1\). The empty edge case has the usual zero residual.

Distinct walks are retained as arrows of the view category. The product along a walk \(q=(i_0,\ldots,i_m)\) is
\[
 T_q=\tau_{i_0i_1}\cdots\tau_{i_{m-1}i_m}.
\]
Concatenation gives the restriction composition law. An exact vertex family defines a coherent walk-view section, and a coherent walk-view section satisfies every edge equation. Thus this graph construction returns to the original ambient realization type rather than creating a disconnected numerical model.

\subsection{Transported residuals and path gains}
\begin{proposition}[The path telescope]\label{KS40}
For a walk \(q=(i_0,\ldots,i_m)\),
\[
 T_qx_{i_m}-x_{i_0}
 =\sum_{k=0}^{m-1}
 (\tau_{i_0i_1}\cdots\tau_{i_{k-1}i_k})(dx)_{i_ki_{k+1}}.
\]
The empty prefix is the identity. If \(\gamma(q)\) is the sum of the products of preceding edge norms, then
\[
 \norm{T_qx_{i_m}-x_{i_0}}\le\gamma(q)\norm{dx}_\infty.
\]
\end{proposition}
\begin{proof}
Expand each summand. Its positive transported endpoint cancels the negative endpoint of the next summand, leaving only the two boundary terms. Apply the triangle inequality and the operator product bound to obtain the estimate.
\end{proof}
\source{KS40}
For isometric maps, \(\gamma(q)\le|q|\). For arbitrary bounded equivalences, route length cannot replace the amplification gain. Reversing a route can have a different gain, even though its transport is the inverse map.

\subsection{Holonomy is the remaining exactness condition}
Write \(T_i=T_{p_i}\). Reconstruction from a proposed root value is
\[
 (R_pz)_i=T_i^{-1}z.
\]
For each directed edge define the root loop operator
\[
 K_{ij}=T_i\tau_{ij}T_j^{-1},\qquad
 K=\bigcap_{(i,j)\in\vec E}\ker(K_{ij}-I).
\]
Each \(K_{ij}\) is the transport around an actual rooted loop. No flatness assumption has been made.

\begin{theorem}[Holonomy classification]\label{KS41}
Evaluation at the root and reconstruction give an equivalence
\[
 K\simeq_L\ker d.
\]
A root value extends to an exact vertex family precisely when it lies in \(K\), and that extension is unique. The space \(K\) is independent of the chosen routes. Reconstruction is route independent on \(K\).
\end{theorem}
\begin{proof}
Multiply the edge defect of a reconstructed family by \(T_i\):
\[
 T_i(dR_pz)_{ij}=K_{ij}z-z.
\]
Thus reconstruction is exact precisely when \(z\in K\). Conversely, an exact family has zero defect along every walk by the telescope, so \(T_ix_i=x_o\). Hence \(x=R_p(x_o)\), which proves uniqueness and both inverse identities.

The root values arising from exact families are independent of the routes used to discover them. This identifies \(K\) intrinsically as the image of root evaluation on \(\ker d\), proving route independence. The same uniqueness proves route independence of reconstruction on \(K\). Boundedness of the finite family of transports and their inverses gives the continuous linear equivalence.
\end{proof}
\source{KS41}
With isometric transports the equivalence is an isometry for the vertex maximum norm, because every reconstructed coordinate has the root norm. With general bounded equivalences, it need not be an isometry.

\plate{09_holonomy}{Tree routes reconstruct from an anchor. The remaining edges close loops and impose fixed-vector equations at that anchor. A nontrivial loop can obstruct an arbitrary seed without preventing every nonzero seed from extending. The fixed space, rather than the identity of the whole loop operator, is the relevant realization condition.}{fig:holonomy}

\subsection{Prescribed observations remain visible}
If a prescribed vertex set \(U\) contains \(o\), a candidate \(x\) has an exact repair preserving \(x|_U\) precisely when
\[
 x_o\in K\quad\text{and}\quad (R_px_o)|_U=x|_U.
\]
Necessity follows from the classification. Sufficiency is the reconstructed family itself. The source identifies this criterion with the earlier obstruction class in \(R/d(\ker p_U)\). A graph seed test is therefore an instance of the same protected correction problem, not a replacement for it.

\chap{Closed range is the stability boundary}\label{sec:closed}
\subsection{Projection exists before a residual estimate}
Let \(E\) be a real Hilbert space and \(R\) a real Banach space. Let \(D:E\to R\) be continuous linear. The kernel \(K=\ker D\) is closed, so its orthogonal projection \(q=P_K\) exists~\cite{mathlib-projection}.

\begin{proposition}[Exact orthogonal repair]\label{KS42}
For every \(x\in E\),
\[
 D(qx)=0,\quad q^2=q,\quad\norm{qx}\le\norm x,
\]
and \(qx\) is the unique closest point to \(x\) in \(K\).
\end{proposition}
\begin{proof}
Orthogonal decomposition gives \(x=qx+(x-qx)\), with the two components in \(K\) and \(K^\perp\). For \(y\in K\),
\[
 \norm{x-y}^2=\norm{x-qx}^2+\norm{qx-y}^2.
\]
This proves best approximation and uniqueness. The other statements follow from the defining projection properties.
\end{proof}
\source{KS42}
No closed range assumption on \(D\) is needed for this projection. Such an assumption enters when the size of the residual is asked to control the distance to the exact space.

\begin{theorem}[Closed range and quantitative exactness]\label{KS43}
The following are equivalent:
\begin{enumerate}
\item \(\ran D\) is closed.
\item There is \(C\ge0\) with \(\norm{x-qx}\le C\norm{Dx}\) for all \(x\).
\item There is \(C\ge0\) such that every \(x\) has some exact \(y\) with \(\norm{x-y}\le C\norm{Dx}\).
\end{enumerate}
When the range is closed, the least admissible constant is the operator norm of the inverse of \(D:K^\perp\to\ran D\). In the zero operator case that inverse has zero dimensional domain and its norm is zero.
\end{theorem}
\begin{proof}
The restriction \(D_\perp:K^\perp\to\ran D\) is bijective. If the range is closed it is Banach, so the bounded inverse theorem makes \(D_\perp^{-1}\) continuous~\cite{mathlib-banach}. Since \(x-qx\in K^\perp\),
\[
 x-qx=D_\perp^{-1}(Dx),
\]
which gives the estimate. Any constant in that estimate bounds the inverse norm by applying it to \(u\in K^\perp\), so the constant is optimal.

Conversely suppose the estimate holds. If \(Du_n\) converges with \(u_n\in K^\perp\), then
\[
 \norm{u_n-u_m}\le C\norm{Du_n-Du_m}.
\]
The sequence \(u_n\) is Cauchy and converges in the closed subspace \(K^\perp\). Continuity of \(D\) puts the residual limit in its range. Hence the range is closed. The second condition implies the third by taking \(y=qx\), while best approximation makes the third imply the second.
\end{proof}
\source{KS43}
The last equivalence is important. A nonlinear choice of exact answers cannot bypass the failure of a uniform residual bound. Failure of closed range is not merely a defect of one selected projection algorithm.

\subsection{A small residual far from exactness}
On \(\ell^2(\N)\), index the standard basis from one and define \(De_n=n^{-1}e_n\). The map is bounded and injective, so its exact space is zero. Yet
\[
 \norm{e_n}=1,\qquad\norm{De_n}=1/n\longrightarrow0.
\]
No finite \(C\) can bound \(\dist(e_n,\ker D)\) by \(C\norm{De_n}\). The range is not closed. For example finite truncations of \((1/n)_{n\ge1}\) belong to the range and converge in \(\ell^2\), but a preimage of the limit would be the constant sequence, which does not belong to \(\ell^2\).

This printed diagonal model illustrates the criterion. The general unit-witness obstruction is present in the recovered Hilbert source. The particular infinite diagonal construction is not represented as a newly compiled Lean theorem in this publication.

\plate{10_closed_range}{Residual size controls distance to the exact space only when transverse directions have a positive lower bound. The diagonal example sends unit vectors to residuals tending to zero, so no uniform inverse gain exists. Exact projection still exists.}{fig:closed}

\subsection{Norm zero versus algebraic zero}
If an additive correction image is not closed, distance to that image can vanish without actual membership. Thus an infimum of residual error can be zero while the exact protected correction is unattained. The algebraic obstruction quotient and its possible seminorm answer different questions. One must identify a closed image, compactness condition, or other attainment principle before replacing one question by the other.

\chap{Holonomy repair and transport amplification}\label{sec:holonomyrepair}
\subsection{An orthogonal anchor, not a global least squares answer}
Return to the graph model with a real Hilbert fiber. The root residual is
\[
 \mathcal H z=(K_{ij}z-z)_{ij},\qquad K=\ker\mathcal H.
\]
Let \(q=P_K\), and define
\[
 Px=R_p(qx_o).
\]
This is always exact, fixes exact families, and is idempotent. It is the unique exact family preserving the feasible component \(qx_o\) of the anchor. Its orthogonality is on the anchor fiber. The vertex and edge norms are still maximum norms.

For example on a two-vertex identity connection, \((0,2)\) with root value zero is repaired to \((0,0)\). The constant family \((1,1)\) is closer in the maximum norm, but it does not preserve the anchor. The operator \(P\) is not being claimed to be a global nearest-point projection in that norm.

\subsection{Every factor in the bound has a meaning}
Define the route gain
\[
 A_p=\max_i\gamma(\overline{p_i}),
\]
using reverse routes so that the error is measured in each destination view. The path telescope gives
\[
 \norm{x-R_p(x_o)}_\infty\le A_p\norm{dx}_\infty.
\]
Let \(N_p\) be the residual normalization \((N_pr)_{ij}=T_ir_{ij}\). The exact identity \(\mathcal H=N_pdR_p\) implies
\[
 \norm{\mathcal H(x_o)}_\infty
 \le\norm{N_p}(\norm d A_p+1)\norm{dx}_\infty.
\]

\begin{theorem}[Amplified holonomy repair]\label{KS44}
If \(\ran\mathcal H\) is closed and \(\kappa\) is its transverse inverse gain, then
\begin{equation}\label{eq:graphgain}
 \norm{x-Px}_\infty\le
 \left[A_p+\norm{R_p}\kappa\norm{N_p}(\norm d A_p+1)\right]
 \norm{dx}_\infty.
\end{equation}
The range condition is necessary for the existence of any uniform bound of distance to an exact family by \(\norm{dx}_\infty\). It is automatic for a fixed finite dimensional fiber, but not uniformly quantitative across all connections.
\end{theorem}
\begin{proof}
Insert the uncorrected anchor reconstruction:
\[
 \norm{x-Px}_\infty
 \le\norm{x-R_px_o}_\infty+
\norm{R_p}\norm{x_o-qx_o}.
\]
Apply the path estimate and Theorem~\ref{KS43} to \(\mathcal H\), then use the normalized residual bound. This proves \eqref{eq:graphgain}.

For necessity, apply a proposed global distance estimate to \(R_pz\), with \(z\in K^\perp\). An exact comparison has root value in \(K\), so root evaluation bounds \(\dist(z,K)\) by the global distance. Since \(dR_pz=N_p^{-1}\mathcal Hz\) and normalization has bounded inverse on the finite residual product, this supplies a finite residual estimate for \(\mathcal H\). Theorem~\ref{KS43} forces closed range. In finite dimension every linear image is closed.
\end{proof}
\source{KS44}
For isometric transports, one can use \(\norm{R_p}\le1\), \(\norm{N_p}\le1\), \(\norm d\le2\), and \(A_p\le r\), where \(r\) is the route radius. The bound reduces to
\[
 \norm{x-Px}_\infty\le[r+\kappa(2r+1)]\norm{dx}_\infty.
\]
The exact repair is route independent for a fixed root and inner product. A chosen upper bound involving route gains need not have the same numerical value for different routes.

\subsection{Flatness does not remove amplification}
Consider a scalar chain with \(n+1\) vertices and coordinates \(e_i(z)=2^iz\). Set \(\tau_{ij}=e_ie_j^{-1}\). This connection is flat, but reverse transport along one edge multiplies norms by two. Let
\[
 x_i=8\cdot2^i-1,\qquad y_i=7\cdot2^i.
\]
The family \(y\) is exact and preserves \(x_o=7\). On a forward edge the residual of \(x\) is \(1/2\), and on its reverse it is \(-1\). Thus
\[
 \norm{dx}_\infty=1,\qquad
 \norm{x-Px}_\infty=2^n-1.
\]
A route length bound of \(n\) fails. This is amplification with trivial holonomy, distinct from poor conditioning of the fixed space of a nearly trivial loop.
\source{KS45}

\subsection{Small holonomy can still be poorly conditioned}
In \(\R^2\), the rational rotation
\[
 U_t=\frac1{1+t^2}
 \begin{pmatrix}1-t^2&-2t\\2t&1-t^2\end{pmatrix}
\]
satisfies
\[
 \norm{(U_t-I)z}^2=\frac{4t^2}{1+t^2}\norm z^2.
\]
For \(t\ne0\) its fixed space is zero and its residual gain is \(\sqrt{1+t^2}/(2|t|)\), which diverges as \(t\to0\). At \(t=0\), the entire fiber is fixed. Therefore small loop discrepancy does not supply a uniform stable realization estimate. The fixed directions and their conditioning must be analyzed together.

\chap{Quantitative descent in a common Hilbert model}\label{sec:quantitative}
\subsection{Local exactness and matching are separate operators}
Let \(E\) be a real Hilbert space of raw local candidates. Let \(D:E\to R\) measure local equation defects and \(C:E\to Y\) measure overlap mismatch, with Banach targets. Put
\[
 L=\ker D,\qquad B=C|_L,
 \qquad M=\ker D\cap\ker C.
\]
The exact matched families form \(M\). First project to \(L\), then project inside \(L\) to \(\ker B\). Denote the resulting map by \(P\).

\begin{theorem}[The two-stage repair is the joint projection]\label{KS46}
The operator \(P\) is the orthogonal projection onto \(M\), equivalently onto the kernel of \(\Delta=(D,C)\). If \(u=x-P_Lx\) and \(v=P_Lx-Px\), then
\[
 u\perp v,\qquad\norm{x-Px}^2=\norm u^2+\norm v^2.
\]
\end{theorem}
\begin{proof}
The output lies in \(L\cap\ker C\). The first error \(u\) belongs to \(L^\perp\), and the second error belongs to \(L\) and is orthogonal to \(M\). Hence their sum is orthogonal to \(M\), which characterizes the joint projection. Their mutual orthogonality gives the identity.
\end{proof}
\source{KS46}
This identity does not require the local and overlap operators to commute. The second projection is taken in the already exact subspace, not in the whole raw space without restriction.

\subsection{A stable repair requires a stable matching problem}
Suppose \(\ran D\) and \(\ran B\) are closed, with residual gains \(a,b\). Then
\[
 \norm u\le a\norm{Dx},\qquad
 \norm v\le b\norm{CP_Lx}
 \le b(\norm{Cx}+\norm C\,a\norm{Dx}).
\]
Therefore
\begin{equation}\label{eq:jointgain}
 \norm{x-Px}\le
 a(1+b\norm C)\norm{Dx}+b\norm{Cx}.
\end{equation}
The Pythagorean identity gives the corresponding squared estimate without adding the two errors before squaring. For the maximum norm on \(R\times Y\), \eqref{eq:jointgain} supplies the sufficient gain \(G=a(1+b\norm C)+b\). It is not asserted to be optimal.

\begin{theorem}[Closed range for the joint defect]\label{KS47}
Assume \(\ran D\) is closed. Then
\[
 \ran(D,C)\text{ is closed}\quad\Longleftrightarrow\quad
 \ran(C|_{\ker D})\text{ is closed}.
\]
Under either condition, a finite joint residual-to-repair bound exists.
\end{theorem}
\begin{proof}
If the joint range is closed and \(Bk_n\to y\) with \(k_n\in\ker D\), the joint values \((0,Bk_n)\) converge to \((0,y)\). Closedness supplies \(k\in\ker D\) with \(Ck=y\), proving the forward implication.

Conversely let \((Dx_n,Cx_n)\to(r,y)\). Write \(x_n=u_n+k_n\), with \(u_n\in(\ker D)^\perp\) and \(k_n\in\ker D\). Closed range of \(D\) implies \(r\in\ran D\) and makes its transverse inverse bounded, so \(u_n\to u\) with \(Du=r\). Then \(Bk_n=Cx_n-Cu_n\to y-Cu\). Closed range of \(B\) supplies \(k\in\ker D\) with \(Ck=y-Cu\). The limit is \((D(u+k),C(u+k))\), so the joint range is closed. The final assertion follows from Theorem~\ref{KS43}, or directly from \eqref{eq:jointgain}.
\end{proof}
\source{KS47}

\plate{11_descent}{Local correction and overlap correction are orthogonal errors when the second problem is restricted to the locally exact subspace. Their composition is the projection onto the joint exact space. The overlap range inside that subspace is the additional stability obligation.}{fig:quantdescent}

This Hilbert model is a quantitative counterpart of gluing, but it is not a norm on an arbitrary categorical stack. It starts from a fixed raw candidate space, fixed norms, and actual local and overlap operators. Changing those choices requires its own comparison.

\chap{Controlled motion of exact spaces}\label{sec:motion}
\subsection{A two-sided condition number theorem}
Let \(D,D':E\to R\) be continuous linear maps from a real Hilbert space to the same normed target. Write \(q=P_{\ker D}\), \(q'=P_{\ker D'}\). Suppose
\[
 \norm{(I-q)x}\le C\norm{Dx},\qquad
 \norm{(I-q')x}\le C'\norm{D'x},\qquad C,C'\ge0,
\]
and \(\norm{D-D'}\le\varepsilon\).

\begin{theorem}[Motion of the exact projector]\label{KS48}
Without either kernel inclusion as a hypothesis,
\begin{equation}\label{eq:motion}
 \norm{q-q'}\le\max(C,C')\varepsilon.
\end{equation}
In particular a uniform residual gain \(M\) makes \(D\mapsto P_{\ker D}\) \(M\)-Lipschitz on the class of operators carrying that gain.
\end{theorem}
\begin{proof}
For \(z\in\ker D'\),
\[
 \norm{(I-q)z}\le C\norm{(D-D')z}\le C\varepsilon\norm z.
\]
Interchanging the roles gives \(\norm{(I-q')w}\le C'\varepsilon\norm w\) on \(\ker D\). Since orthogonal projections are self-adjoint, the latter estimate also bounds \(q(I-q')\) by \(C'\varepsilon\).

For any \(x\), decompose
\[
 (q-q')x=q(I-q')x-(I-q)q'x.
\]
These two terms are orthogonal, the first in \(\ran q\) and the second in its orthogonal complement. Applying the cross estimates to the two orthogonal components \((I-q')x\) and \(q'x\) gives
\[
 \norm{(q-q')x}^2
 \le\varepsilon^2\left((C')^2\norm{(I-q')x}^2+C^2\norm{q'x}^2\right)
 \le\varepsilon^2\max(C,C')^2\norm x^2.
\]
Take square roots and the operator supremum.
\end{proof}
\source{KS48}
There is no smallness threshold when both certificates are already supplied. The estimate can be uninformative when their gains are large. Knowing only the old certificate is a different situation.

For changing input as well as model, the same theorem gives
\[
 \norm{qx-q'y}\le\norm{x-y}+\max(C,C')\varepsilon\norm y.
\]
It also bounds the change in distance to the exact space by \(\max(C,C')\varepsilon\norm x\).

\subsection{Constructing the new certificate}
Assume the target \(R\) is Banach and \(A:R\to E\) is a continuous linear right inverse of \(D\), with \(DA=I\), \(\norm A\le a\), and \(a\varepsilon<1\). Set
\[
 T=(D-D')A,\qquad Z=(I-T)^{-1},\qquad A'=AZ.
\]
The inverse is supplied by the convergent Neumann series because \(\norm T\le a\varepsilon<1\).

\begin{theorem}[Stable surjectivity and derived repair]\label{KS49}
The perturbed map is surjective, with
\[
 D'A'=I,\qquad \norm{A'}\le\frac{a}{1-a\varepsilon}.
\]
Its exact projection satisfies
\[
 \norm{x-q'x}\le\frac{a}{1-a\varepsilon}\norm{D'x},\qquad
 \norm{q-q'}\le\frac{a\varepsilon}{1-a\varepsilon}.
\]
\end{theorem}
\begin{proof}
The identity \(D'A=I-T\) gives \(D'A'=I\), and the Neumann inverse has norm at most \((1-a\varepsilon)^{-1}\). The vector \(x-A'D'x\) is exact for \(D'\), so best approximation bounds the projection error by \(\norm{A'D'x}\). The same argument gives old gain at most \(a\). Apply Theorem~\ref{KS48} with these two gains.
\end{proof}
\source{KS49}
If a measured residual \(y\) satisfies \(\norm{D'x-y}\le\eta\), the distance to the new exact space is at most \(a(\norm y+\eta)/(1-a\varepsilon)\). The uncertainty radius remains a quantity to justify in the application. The theorem does not estimate it statistically.

For a surjective Hilbert-source operator, the right inverse need not be a separate unexplained input. The inverse of \(D:(\ker D)^\perp\to R\), followed by inclusion, constructs it. A known residual certificate bounds its norm. Surjectivity is still substantive, and cannot be assumed for a redundant residual map whose target contains unattainable directions.

\subsection{An equivalence between the moving kernels}
\begin{theorem}[Explicit kernel transport]\label{KS50}
Under the hypotheses of Theorem~\ref{KS49},
\[
 J:\ker D\simeq_L\ker D',\qquad Jx=x-A'D'x,
\]
with inverse \(y\mapsto y-ADy\). Moreover
\[
 \norm{Jx-x}\le\frac{a\varepsilon}{1-a\varepsilon}\norm x.
\]
\end{theorem}
\begin{proof}
Let \(R_0=I-AD\), \(R_1=I-A'D'\). They retract onto their respective kernels. Since \(A'=AZ\), one has \(R_0A'=0\). Also \(D'A=I-T\) gives \(R_1A=0\). Hence \(R_0R_1x=x\) on \(\ker D\), and \(R_1R_0y=y\) on \(\ker D'\). This proves both inverse identities. On \(\ker D\), \(D'x=(D'-D)x\), so the displacement estimate follows from the bound on \(A'\).
\end{proof}
\source{KS50}
The map \(J\) depends on the selected splitting. It is not generally the restriction of the orthogonal projection \(q'\). Equivalence of exact spaces and nearest-point repair answer different questions.

\subsection{An all-real worked model}
Let \(D_t(x,y)=x+ty\). Its kernel is the line \(\{(-tb,b):b\in\R\}\), and its orthogonal projector is
\begin{equation}\label{eq:lineprojector}
 q_t=\frac1{1+t^2}
 \begin{pmatrix}t^2&-t\\-t&1\end{pmatrix}.
\end{equation}
The common right inverse \(r\mapsto(r,0)\) supplies gain at most one for every \(t\), and \(\norm{D_s-D_t}=|s-t|\). The general theorem gives \(\norm{q_s-q_t}\le|s-t|\).

\begin{proposition}[Exact motion of the line]\label{KS51}
For every real \(t\),
\[
 \norm{(q_t-q_0)x}^2=\frac{t^2}{1+t^2}\norm x^2,
 \qquad \norm{q_t-q_0}=\frac{|t|}{\sqrt{1+t^2}}.
\]
In the small perturbation interface, kernel transport sends \((0,b)\) to \((-tb,b)\). Orthogonal projection instead sends it to \((-tb/(1+t^2),b/(1+t^2))\).
\end{proposition}
\begin{proof}
The matrix in \eqref{eq:lineprojector} is symmetric, idempotent, and has range equal to the displayed kernel, so it is the projector. Squaring its symmetric difference from \(q_0\) gives \(t^2/(1+t^2)\) times the identity. This proves the all-vector identity and the operator norm, using a unit vector for equality. The two image formulas follow by substitution. The generic strict-radius interface applies when \(|t|<1\), although the displayed line formulas are valid for every real \(t\).
\end{proof}
\source{KS51}

\plate{12_moving_line}{The exact space can move continuously without remaining equal to the original kernel. Transport along a chosen splitting preserves a coordinate, while orthogonal projection minimizes distance. In the line example these two maps have different formulas even though both land in the same exact space.}{fig:movingline}

\subsection{What fails without the hypotheses}
Take \(D_0=\operatorname{diag}(1,0)\) and \(D_t=\operatorname{diag}(1,t)\). The operator difference has norm \(|t|\). For every nonzero \(t\), however, the new kernel is zero, while the old kernel is the vertical axis. Their projections differ in norm by one. Operator closeness alone does not imply kernel continuity.

More generally, suppose a unit vector \(x\) is old exact but orthogonal to the new kernel. Every new residual gain \(C'\) obeys
\[
 1=\norm{x-q'x}\le C'\norm{D'x}\le C'\varepsilon.
\]
Losing an exact direction under a small operator change forces the new gain to grow. At the strict radius boundary, \(D=I\), \(A=I\), and \(D'=0\) have \(a=\varepsilon=1\), yet \(D'\) is not surjective. Thus the strict condition in Theorem~\ref{KS49} cannot be replaced universally by a non-strict one.
\source{KS52}

\subsection{Motion of the actual descent repair}
Compare pairs \((D,C)\) and \((D',C')\) on the same raw Hilbert space. Suppose each pair has closed local range and closed overlap range on its own locally exact subspace. Let \(G,G'\) be the resulting joint residual gains. If
\[
 \norm{D-D'}\le\varepsilon_D,\qquad
 \norm{C-C'}\le\varepsilon_C,
\]
then the maximum norm on the target product gives
\[
 \norm{(D,C)-(D',C')}\le\max(\varepsilon_D,\varepsilon_C).
\]
Applying Theorem~\ref{KS48} to the joint operators, and using the actual two-stage identity of Theorem~\ref{KS46}, gives
\begin{equation}\label{eq:descentmotion}
 \norm{P_{D,C}-P_{D',C'}}
 \le\max(G,G')\max(\varepsilon_D,\varepsilon_C).
\end{equation}
This is a comparison of the original repair maps, not a replacement definition that happens to have a similar kernel. Both the local equations and overlap conditions can change. The raw space and its inner product remain fixed.
\source{KS53}

\chap{Mathematical meaning and the remaining boundaries}\label{sec:meaning}
\subsection{A relative theory can be complete without being absolute}
The strict space \(\Gamma(F)\) is canonical relative to its diagram: its universal property determines it up to the appropriate unique isomorphism. The additive obstruction is canonical relative to its admissible corrections. A marked realization problem is invariant under the specified comparison equivalences. A repair can be canonical relative to a protected observable, or orthogonal relative to a fixed inner product.

These are meaningful and different kinds of canonicity. They need not culminate in one object independent of every choice of views, coefficients, topology, semantics, and geometry. Some changes of view intentionally forget information. Others strengthen the question. Treating every such change as an equivalence would remove the very distinctions the calculus is designed to preserve.

The appropriate central statement is therefore not that every possible mathematical presentation lives in one fixed universal space. It is that a declared system of representations can be organized into coherent realizations, with explicit extension, comparison, and repair problems. A proposed bridge between two such systems has to show what structure it preserves and what information it loses.

\subsection{Internal realization and intended objects}
Suppose a separate category \(\mathcal S\) contains the intended semantic objects of an application, and a realization construction gives a comparison functor
\[
 \mathcal R:\mathcal S\longrightarrow\mathcal K,
\]
where \(\mathcal K\) is the relevant category or groupoid of coherent data. Defining \(\mathcal K\) does not establish that \(\mathcal R\) is faithful, full, or essentially surjective. These are new mathematical questions about the actual semantic comparison.

Faithfulness asks whether distinct semantic arrows can become indistinguishable. Fullness asks whether every compatible comparison comes from a semantic one. Essential surjectivity asks whether every intended coherent datum is realized by an actual semantic object. Depending on the application, one may seek only a restricted image or a classified obstruction to belonging to it.

This distinction prevents a formal globalization of encoded data from being mistaken for an unrestricted arithmetic or physical existence theorem. The abstract construction does not manufacture missing semantic objects.

\subsection{Every finite test can pass while the infinite problem fails}
\begin{example}[Finite prefixes without a global point]\label{KS54}
Let \(X_n=\{m\in\N:m\ge n\}\), with bonding maps given by inclusion. Every level is nonempty, and every finite prefix has a coherent point. The inverse limit is nevertheless empty.
\end{example}
\begin{proof}
For a prefix ending at \(N\), assign the integer \(N\) at every level. All inclusions preserve it. In a global point, however, the inclusion equations force every coordinate to be the same integer \(m\), while membership at level \(n\) requires \(m\ge n\) for every \(n\). Taking \(n=m+1\) is a contradiction.
\end{proof}
\source{KS54}
This elementary printed counterexample is not added to the recovered Lean library by this edition. It explains why inverse-limit existence needs more than separate finite solvability. Finiteness, compactness, surjectivity, or another coherence principle must be proved in the relevant application. The example does not contradict the strict gluing theorem: a compatible matching family of all prefixes was never supplied.

\subsection{A finite computation that exposes the preservation contract}
Consider a raw vector \(x=(2,0,5)\in\R^3\) with defect \(d(x)=x_1-x_2\). Unconstrained orthogonal repair gives \((1,1,5)\), with squared error two. If the first and third coordinates must be preserved, the unique exact repair is \((2,2,5)\), with squared error four.

Both outputs are exact. Only one satisfies the second preservation problem. Neither theorem contradicts the other because the admissible changes differ. An implementation that reports only zero residual would not distinguish the two tasks.

For a finite linear system, the same idea becomes a direct certificate problem. Let \(A\) encode compatibility and \(P\) the protected values. Choose a raw lift \(x_0\) of the prescribed data and a matrix \(N\) whose columns span \(\ker P\). The extension problem is
\[
 ANz=Ax_0,\qquad y=x_0-Nz.
\]
The class of \(Ax_0\) modulo \(\im(AN)\) is the existence obstruction. The space \(\ker A\cap\ker P\) measures ambiguity. Exact elimination can solve this finite instance, but correctness of an implementation, conditioning in floating point, and the existence of suitable finite encodings are separate obligations.

\subsection{Why this is relevant to machine reasoning}
The framework gives a machine explicit mathematical objects to manipulate: views, restrictions, partial realizations, typed defects, admissible corrections, and comparison witnesses. A successful step can carry both an output and a certificate of what remained fixed. A failed step can return an obstruction in a specified quotient rather than an unexplained failure flag.

The split coordinate theorem is especially direct. In its stated class, the same candidate is represented by protected information and residual data, and repair zeros only the residual coordinate. The descent theorem explains which local certificates can be assembled. The analytic results explain when a small measured defect has a quantitatively reliable meaning.

These observations motivate a machine oriented mathematical language. They do not establish an autonomous reasoning system, an empirical intelligence gain, or a general efficient algorithm. Classical choice, infinite products, quotient membership, and bounded inverse constructions can be mathematically valid without yielding executable finite procedures. Implementation claims need their own specifications and tests.

\subsection{What this monograph closes, and what it leaves open}
The book supplies a coherent account of strict realization spaces, additive and nonabelian extension problems, split tower repair, retained multiscale detail, a genuine composition-site stack, presentation transport, graph holonomy, closed-range stability, and moving exact spaces. Each layer is attached to the data that make it meaningful.

The following broader targets are not conclusions of this edition: arbitrary higher categorical coefficients and infinity-stack descent, unrestricted changes of view doctrine, norm reconstruction for every bounded infinite residual hierarchy, nonlinear or unknown-connection repair, perturbation across changing Hilbert geometries without new comparison data, and semantic essential surjectivity for unrestricted arithmetic or physical applications.

The proof infrastructure also has a separate next task: recover and reconcile any additional historical modules, choose a pinned integrated source lineage, and run a full fresh cumulative compiler and axiom audit on that lineage. This is an integration obligation, not a reason to label the explicit mathematical constructions in this book merely speculative.

\begin{insight}
\textbf{The organizing principle.} Coherent realization, obstruction to extension, freedom of completion, and stability of repair belong to one mathematical conversation. Their unification is useful precisely when it does not identify them with one another.
\end{insight}

\appendix
\chap{Notation and source concordance}\label{app:concordance}
\subsection{Notation}
\begin{longtable}{@{}p{.25\linewidth}p{.7\linewidth}@{}}
\toprule
Symbol & Meaning\\\midrule\endhead
\(\V,F_v,f^*\) & A declared view category, its fibers, and contravariant restrictions.\\
\(\Gamma(F)\) & The strict set of coherent global families.\\
\(\Gamma_U(F),\Ext_U(s)\) & Partial sections on a full patch and their global extension fiber.\\
\(\delta,\partial,d,D\) & Respectively the view differential, tower differential, a model residual, and a Hilbert defect operator. Each target is declared where used.\\
\(p,p_U,p_L\) & A protected observable, restriction to prescribed views, and a splitting dependent innovation observable.\\
\(\mathfrak o_U,H_U\) & The existence obstruction and the group of invisible coherent changes. Additive and nonabelian conventions are stated separately.\\
\(C_j,J_j,Q_N,S_N\) & Coarsener, chosen lift, restored coarse remainder, and retained finite reconstruction.\\
\(c_f,a_v\) & Coherent nonabelian transition labels and gauge values.\\
\(\mathfrak K_G,J_{\rm comp}\) & The cocycle patch groupoid and its composition-support topology.\\
\(\mathcal E_b(R)\) & Realizations equipped with an actual identification of their boundary with \(b\).\\
\(T_i,K_{ij},R_p\) & Transport to the root, the corresponding root loop operator, and reconstruction from the root.\\
\(\mathcal H,q,P\) & Root residual, an orthogonal kernel projector, and the explicitly defined model repair. A graph repair need not be orthogonal in the global maximum norm.\\
\(\simeq,\simeq_L\) & A specified equivalence, and a continuous linear equivalence. Neither notation alone asserts an isometry.\\
\bottomrule
\end{longtable}

\subsection{How to read the correspondence}
The identifiers below are stable editorial identifiers, not counts of new discoveries. A row can combine several source declarations into one exposition. The named declaration, its line range, source hash, and correspondence note are recorded in \code{data/claim_ledger.json} and in the companion Lean guide. Existence of a matching source declaration is not a mechanical proof that every word of an edited statement matches its type. The displayed hypotheses and the printed proof remain part of the review.

\textbf{F} denotes the recovered foundation tree at version 0.20.0. \textbf{M} denotes the separate focused moving-space tree at version 0.32.0. \textbf{S} denotes the separately recovered Hilbert holonomy source. \textbf{W} denotes an additional written example, not a newly compiled Lean declaration. Rows marked F/M refer to different explicitly identified source statements, not an assertion that the two trees have been jointly rebuilt.

\small
\begin{longtable}{@{}p{.09\linewidth}p{.35\linewidth}p{.49\linewidth}@{}}
\toprule ID & Topic & Source modules\\\midrule\endhead
KS01 & Limit universal property & \path{F: ViewSections}\\[3pt]
KS02 & Functorial transport & \path{F: ViewSections}\\[3pt]
KS03 & Terminal evaluation & \path{F: ViewSections}\\[3pt]
KS04 & Strict section gluing & \path{F: SectionGluing}\\[3pt]
KS05 & Object-only cover failure & \path{F: ViewBoundaryExamples}\\[3pt]
KS06 & Correspondence fibers & \path{F: ViewCorrespondence}\\[3pt]
KS07 & Nonempty fibers without a section & \path{F: ViewBoundaryExamples}\\[3pt]
KS08 & Additive view differential & \path{F: SectionObstruction}\\[3pt]
KS09 & Additive extension criterion & \path{F: SectionObstruction}\\[3pt]
KS10 & Extension ambiguity torsor & \path{F: SectionExtensionFibers}\\[3pt]
KS11 & Parity seed example & \path{F: SectionExtensionExample}\\[3pt]
KS12 & Exactness at partial sections & \path{F: SectionExactSequence}\\[3pt]
KS13 & Extension by stages & \path{F: SectionExtensionStages}\\[3pt]
KS14 & Protected correction & \path{F: Correction}\\[3pt]
KS15 & Obstruction witness & \path{F: Correction}\\[3pt]
KS16 & Exact tower maps and realized points & \path{F: Towers}\newline \path{F: Realization}\\[3pt]
KS17 & Tower differential composition & \path{F: Towers}\newline \path{F: TowerCorrection}\\[3pt]
KS18 & Cofinal point equivalence & \path{F: Cofinal}\\[3pt]
KS19 & Geometric operator repair & \path{F: RepairDynamics}\newline \path{F: RepairProjection}\\[3pt]
KS20 & Constructed split tower correction & \path{F: SplitTowerRepair}\\[3pt]
KS21 & Split semantic and residual coordinates & \path{F: SplitTowerCoordinates}\newline \path{F: SplitTowerSemantics}\\[3pt]
KS22 & Growing identity repair & \path{F: GrowingTower}\\[3pt]
KS23 & A bounded residual without bounded primitive & \path{F: IdentityTowerObstruction}\\[3pt]
KS24 & Dependence on target splitting & \path{F: SplittingDependence}\\[3pt]
KS25 & Lossless dyadic residual split & \path{F: DyadicResidualSplit}\\[3pt]
KS26 & Cascade telescope and persistent part & \path{F: CascadeProjections}\newline \path{F: CoarseningCascade}\\[3pt]
KS27 & Recovery by density & \path{F: CascadeDensity}\\[3pt]
KS28 & Dyadic finite-coordinate recovery & \path{F: DyadicInfiniteReconstruction}\\[3pt]
KS29 & Shift cascade and c0 & \path{F: SequenceCascade}\\[3pt]
KS30 & Gauge groupoid and section transport & \path{F: GroupViewSystems}\newline \path{F: CocycleGaugeGroupoid}\\[3pt]
KS31 & Nonabelian extension obstruction & \path{F: NonabelianExtension}\\[3pt]
KS32 & Nonabelian extension torsor & \path{F: NonabelianExtensionFibers}\\[3pt]
KS33 & S3 loop model & \path{F: CocycleExamples}\\[3pt]
KS34 & Effective groupoid descent & \path{F: DescentEquivalence}\newline \path{F: DescentComparison}\\[3pt]
KS35 & Composition topology & \path{F: CompositionSite}\\[3pt]
KS36 & Standard stackhood & \path{F: CompositionStack}\\[3pt]
KS37 & Equivalent view presentations & \path{F: SectionViewEquivalence}\newline \path{F: RelativeSections}\newline \path{F: RelativeAmbiguity}\\[3pt]
KS38 & Identified-boundary realization groupoid & \path{F: MarkedRealization}\\[3pt]
KS39 & Marked transport across a square & \path{F: MarkedTransport}\\[3pt]
KS40 & Transported path telescope & \path{F: AmplifiedTransport}\\[3pt]
KS41 & Holonomy classification & \path{F: AmplifiedHolonomy}\\[3pt]
KS42 & Orthogonal exact repair & \path{M: HilbertKernelProjection}\\[3pt]
KS43 & Closed range equivalence & \path{M: HilbertKernelProjection}\\[3pt]
KS44 & Amplified holonomy gain & \path{F: AmplifiedRepair}\newline \path{S: HilbertHolonomyRepair}\\[3pt]
KS45 & Flat transport amplification & \path{F: AmplifiedExamples}\\[3pt]
KS46 & Two-stage joint projection & \path{M: QuantitativeDescent}\\[3pt]
KS47 & Joint closed range & \path{M: DescentClosedRange}\\[3pt]
KS48 & Two-sided projector motion & \path{M: MovingKernelBounds}\\[3pt]
KS49 & Derived perturbed certificate & \path{M: SplitKernelPerturbation}\\[3pt]
KS50 & Transport between moving kernels & \path{M: SplitKernelPerturbation}\\[3pt]
KS51 & Exact moving-line example & \path{M: MovingKernelExamples}\\[3pt]
KS52 & Rank loss and radius endpoint & \path{M: MovingKernelBounds}\newline \path{M: MovingKernelExamples}\\[3pt]
KS53 & Motion of quantitative descent & \path{M: MovingDescent}\\[3pt]
KS54 & Finite prefixes without a point & W: printed proof only\\[3pt]
\bottomrule\end{longtable}
\normalsize


\par\noindent\begin{minipage}{\linewidth}\subsection{Verification is a separate record}
The foundation archive reports 122 production proof modules, 1,299 named theorem declarations, and 696 indexed constructions. Its recorded strategy includes one complete cumulative source build, followed by a fresh archive replay that reuses a checked 117-module baseline and compiles the five additions and root again. It must not be described as two full clean cumulative builds.

The focused moving-space archive reports 11 production modules and a root, 151 named theorem declarations, and 47 indexed constructions. Its historical replay record reports two complete clean builds of that focused source closure. The archive SHA-256 and the retained source hashes are checked by this publication's integrity script. The focused tree does not replace the whole historical Komorebi Space repository.

The pinned environment in the records is Lean 4.34.0-rc2 with mathlib commit\par\smallskip\noindent\path{aeb72f1e9981b6d35b426b16dca7759be3e7ba06}.\par\smallskip The permitted standard axioms are \code{propext}, \code{Classical.choice}, and \code{Quot.sound}. The compiler and cached dependencies remain within the historical trust boundary. They are not independently reconstructed here.

The current editorial audit records what was actually executed: source and archive integrity, declaration correspondence, exact finite examples, and PDF preflight. It records that a fresh cumulative Lean build was not run. The source archives and their original reports are preserved separately from the newly generated publication audit.

\end{minipage}

\clearpage
\chap{Background and source documents}\label{app:references}
The references identify the standard background and frozen project sources. Printed arguments are supplied for the principal constructions. The dependent-index dyadic reconstruction proof is explicitly an outline, with its complete Lean declaration identified. Live mathlib documentation is a navigation aid, not a replacement for the pinned source revision.

\begin{thebibliography}{9}
\bibitem{stacks-limits} The Stacks Project Authors. \emph{Limits and colimits over preordered sets}. Section 4.21, Tag 002Z. \url{https://stacks.math.columbia.edu/tag/002Z}. Accessed 10 October 2026.
\bibitem{stacks-stack} The Stacks Project Authors. \emph{Stacks}. Section 8.4, Tag 0268. In particular the effective descent and equivalence-of-descent-categories formulations. \url{https://stacks.math.columbia.edu/tag/0268}. Accessed 10 October 2026.
\bibitem{mathlib-stack} The mathlib Community. \emph{CategoryTheory.Sites.Descent.IsStack}. Standard pseudofunctor stack interface. \url{https://leanprover-community.github.io/mathlib4_docs/Mathlib/CategoryTheory/Sites/Descent/IsStack.html}.
\bibitem{mathlib-projection} The mathlib Community. \emph{Analysis.InnerProductSpace.Projection.Basic}. Orthogonal projection onto complete subspaces. \url{https://leanprover-community.github.io/mathlib4_docs/Mathlib/Analysis/InnerProductSpace/Projection/Basic.html}.
\bibitem{mathlib-banach} The mathlib Community. \emph{Analysis.Normed.Operator.Banach}. Banach open mapping and continuous inverse constructions. \url{https://leanprover-community.github.io/mathlib4_docs/Mathlib/Analysis/Normed/Operator/Banach.html}.
\bibitem{foundation} Komorebi P. \emph{Komorebi Space Theory}, version 0.20.0. Recovered source archive and replay evidence, 28 September 2026. Included unchanged under \code{evidence/archives/}. Its source documents include \emph{Ambient Sections and Extension Obstructions}, \emph{Split Tower Coordinates}, \emph{Infinite Reconstruction}, \emph{Nonabelian Coherence}, \emph{Effective Groupoid Descent}, \emph{A Stack on the Composition Site}, and \emph{Nonisometric Transport and Amplified Reconstruction}.
\bibitem{moving} Komorebi P. \emph{Komorebi Space: Controlled Motion of Exact Realization Spaces}, focused version 0.32.0. Recovered source archive and fresh replay evidence, 30 September 2026. Included unchanged under \code{evidence/archives/}.
\bibitem{hilbert-extra} Komorebi P. \emph{HilbertHolonomyRepair.lean}. Separately recovered source implementing closed-range amplified holonomy reconstruction. Included under \code{evidence/selected/}. It is identified separately rather than assigned the verification coverage of a different archived source closure.
\end{thebibliography}

\subsection*{Editorial source policy}
Earlier source documents contain historical scope statements and counts. Later checked developments can supersede an earlier statement that a topic is a target. In particular the old outline's stack target is not the status of the later standard stack theorem. This monograph uses the actual relevant declarations rather than treating the earliest outline as a current completion ledger.



\end{document}
